% ---------------------------------------------------------------------
% CHAPTER 1 — INTRODUCTION
% Job: establish the problem and the single thesis. NO overarching RQs here.
% ---------------------------------------------------------------------
\chapter{Introduction \label{ch:introduction}}

% 1.1 The problem, stated once, concretely.
%   - Architecture drifts/erodes; documentation goes stale; mappings are missing.
%   - SAR and C2A mapping defined and distinguished (recovery = discover the
%     decomposition; mapping = assign entities given partial known architecture).
%   - Cost anchor: expert C2A mapping is ~107 person-hours/system; Chromium took
%     ~2 years. This is the "why automate" evidence, not rhetoric.
\section[Architectural Erosion and Recovery]{Software Architecture Erosion and the Recovery}

% 1.2 Why the obvious approaches are insufficient (motivates the whole thesis).
%   - Dependency-only clustering: "connected by an edge" != "same module".
%   - Hand-tuned fusion (SARIF): fixed weights don't generalize across systems/scale.
%   - The gap: no single signal suffices, and hand-weighting them doesn't either.
\section{Why architectural signal is hard to use}

% 1.3 The thesis statement (the spine). State it plainly:
%   Architectural structure is distributed across several implementation signals
%   (typed dependencies, folder organization, lexical identifiers, member-level
%   detail). Learned graph representations integrate them better than hand-tuned
%   fusion; and recovery is bounded not by the encoder but by whether the target
%   architecture is reflected in the chosen signal and whether the model's reach
%   matches system scale and granularity.
\section{Thesis statement and scope}

% 1.4 Contributions of the thesis as a whole (not per-paper yet).
%   - C1: learned representations beat hand-tuned fusion, most clearly at scale.
%   - C2: isolation of what each signal contributes and where it fails.
%   - C3: the same signals/systems hold across unsupervised recovery AND
%         supervised mapping (the duality that no single paper can claim).
%   - C4: directory structure as a strong architectural proxy, shown by removing it.
\section{Contributions}

% 1.5 Included papers + author roles + a one-line "what each isolates".
\section{Included publications}

% 1.6 Thesis outline (short).
\section{Outline}
% \section{Motivation: Architecture Erosion and the Cost of Technical Debt}
% \label{sec:motivation}

% As software systems evolve through refactoring, feature additions, and ongoing maintenance, the implemented system gradually diverges from its intended architectural design. This architectural drift, or erosion, degrades maintainability, understandability, and long-term evolvability~\cite{bass2012software,de2012controlling}. Static Architecture Conformance Checking (SACC) techniques, such as reflexion modeling, address this problem by comparing implementation-level dependencies against architecture-level relations to identify convergent, divergent, and absent dependencies~\cite{reflexion1, SACCcomparison}. Applying these techniques presupposes that a decomposition of the system into architectural modules is available. In many systems, no such decomposition is documented or maintained, and it must instead be recovered directly from the implementation.
 
% Software Architecture Recovery (SAR) addresses this first problem: inferring a system's modular decomposition from its source code, without assuming that an intended architecture is already known. Recovering an architecture manually is labor-intensive and does not scale; recovering the decomposition of a medium-sized system can require hundreds of hours of expert effort~\cite{garcia2013}. This motivates automated recovery techniques that infer modular structure directly from implementation artifacts such as dependencies, folder organization, and code identifiers.
 
% Recovery, however, is not the only setting in which this problem arises. In many systems, an intended architecture already exists, but the mapping between implementation-level entities and architectural modules is incomplete, outdated, or was never fully documented~\cite{ali2018architecture}. Here, the task is not to discover a decomposition from scratch, but to propagate a small amount of known mapping to the remainder of the system, a problem generally referred to as code-to-architecture (C2A) mapping. Both problems, recovery and mapping, require inferring architectural structure from the same underlying implementation signals, differing primarily in how much architectural knowledge is available as supervision: none, in recovery, and partial, in mapping. Maintaining this mapping is a growing concern as the volume of machine-generated code entering software systems increases, since such code is produced without visibility into a system's module structure~\cite{ehsani2025fastercode,wu2025clones}.
 
% Existing techniques for both problems typically combine multiple implementation signals, such as static dependencies, folder organization, and lexical identifiers, through hand-crafted similarity functions and manually tuned weights~\cite{sarif,sade}. Because dependency patterns and architectural styles differ across systems, such fixed heuristics do not generalize consistently. This motivates a data-driven alternative: learning, rather than hand-specifying, how these signals should be combined and weighted for a given system.

% Every non-trivial software system begins with an intended architecture: a decomposition into modules, layers, or subsystems that encodes the designers' understanding of how responsibilities should be separated and how components should interact. As a system evolves under the pressure of deadlines, changing requirements, and developer turnover, its \emph{implemented} architecture tends to drift away from this intended structure. This phenomenon, known as \textbf{architecture erosion} (or architectural degradation), has been studied extensively in the software engineering literature and is consistently linked to rising maintenance costs, reduced developer productivity, and increased defect rates as systems age~\citep{li2022erosion}.

% The financial framing for this problem was introduced by Ward Cunningham, who coined the \textbf{technical debt} metaphor to describe the trade-off between short-term development speed and long-term maintainability:

% \textit{``Shipping first time code is like going into debt\ldots{} the danger occurs when the debt is not repaid.''}

% \hfill --- \citet{cunningham1992wycash}

% When this debt accumulates specifically at the level of a system's structure, in its module boundaries, layering, and dependency organization, it is termed \textbf{architectural technical debt (ATD)}: design or implementation constructs that are expedient in the short term but make future changes more costly, precisely because they violate or obscure the system's intended decomposition~\citep{kruchten2012technicaldebt,nord2012metric}. The technical debt community formalizes this short-term/long-term trade-off using two further concepts borrowed directly from finance: the \textbf{principal}, the cost of the rework required to bring a system back to a well-architected state, and the \textbf{interest}, the recurring extra cost paid on every subsequent change until that rework happens~\citep{verdecchia2020grounded}. Code-level debt is often locally visible as a "smell" in a single file and can be readily flagged by static analyzers. Architectural debt is different: it lives in the relationships \emph{between} components, is invisible to any single developer working within one module, and tends to be discovered only once its accumulated interest, the extra effort required to make a change that should have been simple, becomes large enough to notice. Existing attempts to make ATD measurable, such as the ATDx index, aggregate violations of a manually curated set of architectural rules into a single score~\citep{verdecchia2020atdx}. That approach is useful, but only as good as the rules an analyst thought to specify in advance. This asymmetry, between the difficulty of \emph{detecting} architectural debt (diffuse, relational, and only partially captured by hand-picked rules) and the ease of \emph{incurring} it, motivates the two complementary problems this dissertation addresses.

% The first problem is a matter of \textbf{diagnosis}: once documentation has gone stale or was never created, how can we recover an accurate, up-to-date picture of a system's actual architecture directly from its code, so that architectural technical debt, and specifically its accumulated \emph{principal}, can even be identified in the first place? The second problem is a matter of \textbf{prevention}: given that a system already has a known, intended decomposition into modules, how can new or previously unclassified code be correctly and consistently placed within it, so that architectural debt stops accumulating and stops charging \emph{interest} on every subsequent change, simply because developers lack automated guidance on where new code belongs? Both problems are, at their core, about recovering or preserving the correspondence between a codebase's structure and its architecture: one after the correspondence has been lost and debt has already accrued, the other before it can be lost, preventing debt from accruing at all. This dissertation studies both, using graph neural network-based representation learning as the common technical vehicle, and treats technical debt as the motivation revisited throughout every subsequent chapter.

% \section{Problem Statement: Two Complementary Architectural Tasks}
% \label{sec:problem-statement}

% We study two related but distinct tasks throughout this body of work.

% \paragraph{Software Architecture Recovery (SAR)} is the \emph{unsupervised} problem of inferring a module-level decomposition of a software system directly from its source code, with no ground-truth structure provided as input. The output is typically a clustering of code entities (files, classes) into architecturally meaningful groups, evaluated against an expert-defined or historically documented "true" decomposition. SAR is the appropriate task when documentation is absent, outdated, or untrustworthy: the diagnostic scenario above.

% \paragraph{Software/Code-to-Architecture Mapping (SC2AM)} is the \emph{semi-supervised} problem of assigning code entities to an existing, partially known architectural decomposition. Here, a subset of entities already carries a trusted module label, from documentation, prior recovery, or developer annotation, and the task is to propagate this labeling to the remaining, unassigned entities. SC2AM is the appropriate task when an intended architecture is known and the goal is to keep newly added or previously unclassified code consistent with it: the preventive scenario above.

% Read through the technical debt lens established in Section~\ref{sec:motivation}, these two tasks map directly onto the diagnosis/prevention split. SAR estimates a system's currently accumulated architectural debt principal, how far the implemented structure has drifted from a defensible, well-architected baseline. SC2AM works to prevent that principal from growing further, by keeping every newly added or previously unclassified code entity aligned with the intended decomposition and thereby avoiding the interest that would otherwise be paid on every subsequent change built on a misplaced entity. Every paper in this dissertation can be read as offering a more accurate or more scalable instrument for one of these two debt-management functions.

% The distinction in supervision, none for SAR, partial for SC2AM, drives substantially different modeling choices throughout this dissertation, clustering versus classification-style propagation respectively, even though both tasks operate on the same underlying representation of a codebase as a graph of interrelated entities. This dissertation organizes around that split into two threads. Thread A, on code-to-architecture mapping, comprises three works of increasing sophistication: a foundational graph convolutional approach, an extension that isolates the graph's actual contribution against textual and graph-free baselines, and a further extension that applies a heterogeneous graph transformer over folder- and member-augmented graphs. Thread B, on architecture recovery, comprises three works exploring, in turn, contrastive heterogeneous graph attention, graph auto-encoding with combined structural and textual signal, and a systematic comparison of dependency-only versus folder-augmented graph representations.

% \section{Why Graph Neural Networks}
% \label{sec:why-gnns}

% A codebase is not naturally a table of independent feature vectors. It is a network of interdependent entities. Files import and reference one another, classes inherit and compose, methods call one another across module boundaries, and this whole structure sits inside a folder hierarchy that itself encodes organizational intent. Because both SAR and SC2AM are fundamentally about recovering or preserving groupings over this network of entities, methods that can directly exploit relational structure have a natural advantage over methods that treat each entity independently.

% Prior approaches to both problems have largely relied on heuristic clustering over hand-engineered similarity metrics, or information-retrieval-style techniques that treat source files primarily as bags of textual tokens, such as identifier or comment similarity, with structural information folded in, if at all, as an auxiliary feature rather than as the primary learning signal. Such approaches can miss higher-order structural patterns: two files that never directly depend on each other but sit in the same tightly-coupled neighborhood of the dependency graph are architecturally related in a way that a purely local or textual similarity measure will not capture.

% Graph Neural Networks (GNNs) address this gap directly. They operate over the dependency, folder, and member-interaction relationships of a codebase as first-class structure, propagating and aggregating information along these relationships to produce representations of code entities that reflect both their local properties and their position within the broader architectural network. This has a direct bearing on the technical-debt-measurement problem raised in Section~\ref{sec:motivation}: indices such as ATDx are only as complete as the architectural rules an analyst manually specifies in advance, and can therefore miss debt patterns that were never anticipated. A learned graph representation does not require debt-indicating patterns to be enumerated up front. It is trained to reconstruct or exploit the actual structure of a system, and so can surface architecturally meaningful groupings, and by extension deviations from them, that a fixed rule set was never designed to check for. This dissertation's body of work traces a specific and deliberate evolution in how this idea is operationalized, across the type of graph built from the codebase, the GNN encoder architecture applied to it, and the additional input signals fused into it, which Chapter~\ref{ch:background} introduces technically and Chapters~\ref{ch:thread-a}--\ref{ch:thread-b} trace historically across the two threads.

% \section{Research Questions}
% \label{sec:research-questions}

% The papers making up this dissertation, read individually, each pose a locally-scoped research question about one representation, one encoder, or one signal. Read together, they trace answers to one overarching question:

% \begin{quote}
% \textbf{Overarching RQ.} How can graph neural network-based representation learning be effectively designed and adapted to support both software architecture recovery (unsupervised) and code-to-architecture mapping (semi-supervised), and what governs the effectiveness of different graph representations and input signals across these two tasks?
% \end{quote}

% Three sub-questions unpack this claim, each answered with evidence drawn from multiple papers rather than any single one. The first concerns representation: what graph representation of a codebase, whether a single-relation dependency graph, a typed or heterogeneous dependency graph, one augmented with folder hierarchy, or one further augmented with member-level (method/field) interactions, best supports architecture-related learning, and does the answer differ between recovery and mapping? Read through the debt lens, this question asks which relations in a codebase actually carry debt-relevant signal: which structural cues correlate with where implemented structure has drifted from intended structure. The second concerns encoder architecture: which GNN encoder families, among graph convolutional and graph attention networks, graph auto-encoders trained with contrastive learning, and heterogeneous graph transformers, most effectively exploit these representations for each task? Read through the debt lens, this question asks which encoder best exposes the latent groupings and misassignments that indicate accumulated principal, without requiring the hand-specified rules that indices like ATDx depend on. The third concerns signal attribution and cost: under what conditions, including system scale, the structural cohesion or homophily of the dependency graph, and the availability of labels, do additional signals such as textual semantics, folder structure, and member-level edges yield meaningful improvements over structure alone, and what costs in data engineering effort, computational cost, and hyperparameter sensitivity do they introduce? This third question is itself a debt trade-off: every additional signal amounts to interest paid up front, in engineering and compute cost, against a hoped-for reduction in the principal, meaning better recovery or mapping accuracy. It asks when that trade is worth making.

% This framing is not imposed retrospectively. The most recent published work in Thread B poses almost exactly this question in its own motivation, asking how far graph representation learning can go from dependencies alone, and when additional signals are needed. This dissertation treats that question as the connective thread running through the entire body of work, on both the recovery and mapping sides.

% \section{Overview of the Research Trajectory}
% \label{sec:trajectory-overview}

% \textbf{Thread A, Code-to-Architecture Mapping,} begins with a foundational comparison of homogeneous graph convolutional networks against relational (heterogeneous) graph convolutional networks for semi-supervised mapping, combined with iterative self-training to propagate labels with dynamic confidence thresholding, evaluated on eight Java systems. It is then extended by isolating the graph's actual contribution: introducing a graph attention encoder and a graph-free baseline to determine how much of the method's effectiveness is attributable to structure versus text, and scaling evaluation to ten systems including an industrial-scale, eleven-million-line C/C++ codebase. The thread's most advanced work further enriches the graph itself, adding folder and member-level nodes to the previously file-only dependency graph, and applies a heterogeneous graph transformer with type-aware attention, achieving strong results even at very low labeling rates.

% \textbf{Thread B, Software Architecture Recovery,} begins with a heterogeneous graph attention encoder trained via contrastive Deep Graph Infomax learning, optionally incorporating edge-level attributes describing member-interaction type and frequency, evaluated on a small set of four systems. A subsequent work reframes the problem around a graph auto-encoder, combining dependency structure, directory information, and Word2Vec-derived identifier text into a richer node representation, evaluated at industrial scale, including Chromium, against a wide set of established recovery baselines. The thread's most recent work then deliberately narrows scope to isolate one specific question left open by that combination: how much of the auto-encoder's success comes from dependency structure alone, versus from folder augmentation, versus from the added textual signal. It does so by systematically comparing dependency-only and folder-augmented graph representations across eleven systems while holding textual features constant.

% Across both threads, one empirical pattern recurs independently: the effectiveness of dependency-graph-based methods is strongly tied to a system's \textbf{structural cohesion (homophily)}, the extent to which dependencies stay within, rather than cross, module boundaries. Low homophily is, in effect, a direct symptom of architectural technical debt. A dependency that crosses an intended module boundary is precisely the kind of structural violation practitioners report as a recognizable symptom of accumulated ATD in empirical studies of how architects experience and describe debt in practice~\citep{verdecchia2020grounded}. Folder and organizational signal help most precisely where this cohesion is weak, which is precisely where debt has already accumulated most heavily, and dependency structure alone has become the least reliable diagnostic signal. This finding recurs across the recovery and mapping threads and is treated in this dissertation as the closest thing to a unifying scientific claim across the whole body of work: the papers collected here are, together, tracking how the reliability of graph-based debt diagnosis and prevention degrades as debt itself increases, and what additional signal is needed to compensate.

% \section{Contributions}
% \label{sec:contributions}

% At the level of the dissertation as a whole, rather than any individual paper, this body of work makes five contributions. It gives a systematic, multi-paper account of how graph representation choice, whether single-relation, typed, folder-augmented, or member-augmented, interacts with encoder architecture, whether convolutional, attentional, auto-encoding, or transformer-based, across two related but distinct architecture-related learning tasks. It provides empirical evidence, replicated across independent studies and benchmark systems, that dependency-graph homophily is a key moderating factor for when structural signal alone suffices and when it must be augmented with organizational or textual signal, in effect a moderating factor for how much architectural technical debt a system already carries. It demonstrates, across systems ranging from small open-source Java projects to an eleven-million-line industrial C/C++ codebase, that these methods scale to realistic system sizes without prohibitive computational cost, making them viable as continuous, low-overhead debt-monitoring instruments rather than one-off audits. It traces a comparative architectural progression, from graph convolutional networks through graph attention and auto-encoding to heterogeneous graph transformers, applied consistently to both the recovery and mapping problems, showing how advances developed for one task inform the other. And it offers a learned alternative to hand-crafted architectural technical debt indices such as ATDx, which aggregate violations of manually curated architectural rules into a single score~\citep{verdecchia2020atdx}: rather than requiring an analyst to enumerate debt-indicating rules in advance, the representations developed across this dissertation surface debt-relevant structure directly from dependency, folder, and member-level relationships, and can be evaluated against such indices in future work.

% \section{Structure of the Dissertation}
% \label{sec:structure}

% % The remainder of this kappa is organized as follows. Chapter~\ref{ch:background} provides the technical background necessary to read the two trajectory chapters without prior familiarity with the individual papers: it defines SAR and SC2AM formally, and introduces the GNN architectures used across the body of work. Chapter~\ref{ch:thread-a} traces the research trajectory of Thread A (code-to-architecture mapping). Chapter~\ref{ch:thread-b} traces the research trajectory of Thread B (architecture recovery). Chapter~\ref{ch:synthesis} synthesizes findings across both threads, comparing architectures, representations, and signals directly, and presenting the homophily finding and its implications. Chapter~\ref{ch:limitations} discusses limitations and threats to validity, synthesized across the whole body of work rather than repeated per paper, including, for each paper, which forms of technical debt its methods can and cannot make visible, closing the diagnosis/prevention framing introduced in this chapter. Chapter~\ref{ch:conclusion} concludes with the dissertation's contributions and directions for future work.